% =========================
% Chapter 8
% =========================
\chapter{Conclusion and Future Work}
\label{ch:conclusion}

\section{Summary}

This project set out to answer whether assistive guidance for blind users can be
reliable and context-aware while minimizing cognitive load --- and answered it by
building Lumen: a task-scoped assistive system delivered entirely through a phone
browser and a single Python server. Perception activates only inside explicit,
voice-initiated tasks governed by a finite state machine; guidance is structured,
disciplined, and bounded; and when the task ends, the system is silent.

Three task families were delivered. Object Allocation finds requested objects and
guides the user to them with per-class distance reasoning; Reach Guidance steers the
user's hand to the target and recognizes the grasp itself; and Navigation implements
goal-directed exploration --- to our knowledge an unusual point in the assistive
design space: no maps, no beacons, no SLAM, but a compass-anchored room scan,
semantic arrival detection, a safety-hardened door perception funnel built around a
custom-trained detector (mAP@50 $=0.946$), doorway-transit detection that treats the
camera as an odometer, and an obstacle watchdog fusing named objects with
class-agnostic floor segmentation.

The engineering approach proved as important as the algorithms. Isolating all
decision logic from models and I/O made the system's behavior verifiable --- 286
automated tests, including complete scripted navigation journeys, run in under a
second --- while live trials in real apartments drove the perception design through
a dozen documented failure-mechanism cycles. The result is a system whose
safety-relevant behaviors are not aspirations but tested properties.

\section{Future Work}

\begin{itemize}[nosep]
  \item \textbf{Conversational direction hints (M3/M4).} Ask ``I see two doors ---
        which way is the kitchen?'' and parse left/right/behind answers; remember
        entered-from doors, add breadcrumb backtracking, and prefer unvisited doors
        --- turning the room cap from a check-in into genuine exploration memory.
  \item \textbf{Hard-negative door fine-tune.} Collect the wall footage that fools
        all three models and fine-tune the detector on it --- the durable fix the
        funnel currently substitutes for.
  \item \textbf{A recorded-clip evaluation harness.} Replay annotated phone clips
        through the full perception stack for regression-safe tuning and
        quantitative per-layer ablations (funnel with/without the verifier;
        floor segmentation vs.\ depth vs.\ named-objects-only).
  \item \textbf{A formal user study with blind participants.} Task success, time,
        and standardized workload/usability instruments --- the evaluation that
        sighted-tester trials cannot substitute for, conducted with appropriate
        safety protocols.
  \item \textbf{Offline speech synthesis.} A local TTS fallback to remove the last
        online dependency.
  \item \textbf{Richer vocabularies and languages.} More destination room
        types, Arabic voice interaction, and open-vocabulary object finding.
\end{itemize}

\section{Closing Remark}

The deepest lesson of this project is architectural: treating the user not as a
consumer of narration but as a collaborator in a bounded task --- contributing the
coarse spatial knowledge only they have, while the system contributes the precise
perception only it has --- yields assistance that is simultaneously more capable and
less burdensome. We believe this human-in-the-loop, task-scoped pattern generalizes
well beyond the three tasks built here.
